% GENERATED FILE. Edit canonical lessons, then run npm run generate:books.
% canonical-source-hash: fnv1a64:b5e606caabf45036
% canonical-lessons: ZH-C64-duibuqi, ZH-W64-qi-rise, ZH-C64-bang, ZH-C64-bangmang, ZH-C64-jie, ZH-C64-huan

\chapter{Asking for Help}
\label{ch:zh-asking-for-help}
\hlchaptermodality{64}

\begin{chapteropening}
\textbf{By the end of this chapter:} \emph{I can say sorry, ask for help, borrow something and give it back.}

Complete the last lesson of chapter 64: I can say sorry, ask for help, borrow something and give it back.
\end{chapteropening}

\section[\texorpdfstring{\zh{对不起} (duìbuqǐ)}{对不起 (duìbuqǐ)}]{\zh{对不起} (duìbuqǐ) — sorry}

\label{lesson:ZH-C64-duibuqi}



\begin{quote}
Before the new one: say the Mandarin for identifying number, then the Mandarin for to look for.
\end{quote}

\subsection*{You'll want to know: \zh{对不起}}
\textbf{\zh{对不起}} — \emph{duìbuqǐ} — \textquotedblleft{}sorry\textquotedblright{}.

\begin{quote}
\emph{duìbuqǐ wǒ bù dǒng nǐ de huà} — Sorry, I don't understand what you're saying.
\end{quote}

The new character, in pieces: \zh{起} — the walking piece read zǒu along the left and the bottom, carrying a three-stroke piece read jǐ, 'self', on its back: not the jǐ of \zh{机}. The next lesson puts it in your hand.

Its middle syllable loses its tone: duìbuqǐ, with bu light between two full syllables.

\begin{cousinweb}
\textbf{\zh{对不起}} (\emph{duìbuqǐ}) is built from \zh{对} 'to answer, to match', \zh{不} 'not' and \zh{起} 'to rise', all inherited from Middle Chinese, \emph{twojH}, \emph{pjut} and \emph{khiX} (Guangyun, in Baxter's transcription). 'Not able to face you' is the usual explanation, but it is not settled.
\end{cousinweb}

\begin{grammarlens}[title={What you've built}]
One more word for asking politely.
\end{grammarlens}

\subsection*{Your turn}
\begin{itemize}
\raggedright
  \item \emph{Say it:} \emph{duìbuqǐ}
  \item \emph{Say it:} \emph{duìbuqǐ}, once more
  \item \emph{Say it:} \emph{duìbuqǐ wǒ bù dǒng nǐ de huà} — Sorry, I don't understand what you're saying.
  \item \emph{Recall:} say the Mandarin for identifying number, then the Mandarin for to look for, then say \emph{duìbuqǐ} again
\end{itemize}

\begin{tcolorbox}[breakable,colback=teal!4,colframe=teal!35!black,title={Before you move on}]
What does \zh{对不起} mean? (\textquotedblleft{}Sorry\textquotedblright{}.) Say it once more.
\end{tcolorbox}
\section[\texorpdfstring{\zh{起} (qǐ)}{起 (qǐ)}]{\zh{起} — the walking piece, carrying jǐ}

\label{lesson:ZH-W64-qi-rise}



\begin{quote}
The word first, then the character inside it.

\begin{itemize}
\raggedright
  \item \emph{Recall:} say \emph{duìbuqǐ} — \textbf{R1}, one lesson back, before it decays
\end{itemize}

\begin{quote}
\zh{对} \zh{不}
\end{quote}

\emph{Read it:} each one aloud, with nothing said first — \emph{duì}, \emph{bù}
\end{quote}

\subsection*{Script — the walking piece, carrying jǐ}
\hlblockfigure{\detokenize{figures/ZH-W64-qi-rise-filmstrip.pdf}}{How \zh{起} is written, stroke by stroke}

\begin{quote}
\zh{起}
\end{quote}

\textbf{Two pieces: the walking piece first, then the piece on its back.}

\textbf{The walking piece, seven strokes, read zǒu}, as in the word for to walk: a short horizontal, a vertical down through it, a longer horizontal; then below, a short vertical with a short bar to its right, a left-falling stroke, and last a long stroke that sweeps right along the whole bottom.

\textbf{The right piece, three strokes, read jǐ, 'self'}, not the jǐ of \zh{机}, sitting on the sweep: a horizontal that turns down, a short horizontal, then a vertical that bends right along the bottom and hooks up.

7 + 3 = ten.

\subsection*{Your turn}
\begin{itemize}
\raggedright
  \item \emph{Write it:} \zh{起}, all ten, with the model in front of you
  \item \emph{Cover:} write it again from memory, and count the pieces, not the strokes
\end{itemize}

\begin{tcolorbox}[breakable,colback=teal!4,colframe=teal!35!black,title={Before you move on}]
How many pieces in \zh{起}? (\textbf{Two} — the walking piece and jǐ.) Which stroke runs under the whole character? (\textbf{The long sweep, the seventh}.) Which character comes before it in \zh{对不起}? (\textbf{\zh{不}}.)

Source: \href{https://raw.githubusercontent.com/chanind/hanzi-writer-data/68d10a4b21150cae5e1ebbd223eed289cf32d90c/data/\%E8\%B5\%B7.json}{Hanzi Writer Data \zh{起}, pinned stroke paths and medians}.
\end{tcolorbox}
\section[bāng]{bāng — to help}

\label{lesson:ZH-C64-bang}



\begin{quote}
Before the new one: say the Mandarin for to look for, then the Mandarin for sorry.
\end{quote}

\subsection*{You'll want to know: bāng}
\emph{bāng} — \textquotedblleft{}to help\textquotedblright{}.

\begin{quote}
\emph{nǐ bāng wǒ xiě hǎo bu hǎo} — Will you help me write it?
\end{quote}

Its characters come later; for now it is a word for the ear.

\begin{cousinweb}
\emph{bāng} is inherited from Middle Chinese \emph{pang} (Guangyun, in Baxter's transcription), glossed 'to mend shoes'. The sense 'to help' is not given there, so how it got it is not settled.
\end{cousinweb}

\begin{grammarlens}[title={What you've built}]
One more word for asking politely.
\end{grammarlens}

\subsection*{Your turn}
\begin{itemize}
\raggedright
  \item \emph{Say it:} \emph{bāng}
  \item \emph{Say it:} \emph{bāng}, once more
  \item \emph{Say it:} \emph{nǐ bāng wǒ xiě hǎo bu hǎo} — Will you help me write it?
  \item \emph{Recall:} say the Mandarin for to look for, then the Mandarin for sorry, then say \emph{bāng} again
\end{itemize}

\begin{tcolorbox}[breakable,colback=teal!4,colframe=teal!35!black,title={Before you move on}]
What does bāng mean? (\textquotedblleft{}To help\textquotedblright{}.) Say it once more.
\end{tcolorbox}
\section[bāngmáng]{bāngmáng — to help out; to do a favour}

\label{lesson:ZH-C64-bangmang}



\begin{quote}
Before the new one: say the Mandarin for sorry, then the Mandarin for to help.
\end{quote}

\subsection*{You'll want to know: bāngmáng}
\emph{bāngmáng} — \textquotedblleft{}to help out; to do a favour\textquotedblright{}.

\begin{quote}
\emph{wǒ xiǎng bāngmáng} — I'd like to help.
\end{quote}

Its characters come later; for now it is a word for the ear.

\begin{cousinweb}
\emph{bāngmáng} is built from bāng 'to help' and máng 'busy', both inherited from Middle Chinese, \emph{pang} and \emph{mang} (Guangyun, in Baxter's transcription). The dictionary glosses \emph{mang} as 'afraid', so the sense busy is not given there; helping with what keeps somebody busy is the likely picture, but it is not settled.
\end{cousinweb}

\begin{grammarlens}[title={What you've built}]
One more word for asking politely.
\end{grammarlens}

\subsection*{Your turn}
\begin{itemize}
\raggedright
  \item \emph{Say it:} \emph{bāngmáng}
  \item \emph{Say it:} \emph{bāngmáng}, once more
  \item \emph{Say it:} \emph{wǒ xiǎng bāngmáng} — I'd like to help.
  \item \emph{Recall:} say the Mandarin for sorry, then the Mandarin for to help, then say \emph{bāngmáng} again
\end{itemize}

\begin{tcolorbox}[breakable,colback=teal!4,colframe=teal!35!black,title={Before you move on}]
What does bāngmáng mean? (\textquotedblleft{}To help out; to do a favour\textquotedblright{}.) Say it once more.
\end{tcolorbox}
\section[jiè]{jiè — to borrow; to lend}

\label{lesson:ZH-C64-jie}



\begin{quote}
Before the new one: say the Mandarin for to help, then the Mandarin for to do a favour.
\end{quote}

\subsection*{You'll want to know: jiè}
\emph{jiè} — \textquotedblleft{}to borrow; to lend\textquotedblright{}.

\begin{quote}
\emph{wǒ jiè le tā de shū} — I borrowed his book.
\end{quote}

Its characters come later; for now it is a word for the ear.

\begin{cousinweb}
\emph{jiè} is inherited from Middle Chinese \emph{tsjaeH} (Guangyun, in Baxter's transcription), glossed 'to borrow'; today the same word also lends.
\end{cousinweb}

\begin{grammarlens}[title={What you've built}]
One more word for asking politely.
\end{grammarlens}

\subsection*{Your turn}
\begin{itemize}
\raggedright
  \item \emph{Say it:} \emph{jiè}
  \item \emph{Say it:} \emph{jiè}, once more
  \item \emph{Say it:} \emph{wǒ jiè le tā de shū} — I borrowed his book.
  \item \emph{Recall:} say the Mandarin for to help, then the Mandarin for to do a favour, then say \emph{jiè} again
\end{itemize}

\begin{tcolorbox}[breakable,colback=teal!4,colframe=teal!35!black,title={Before you move on}]
What does jiè mean? (\textquotedblleft{}To borrow; to lend\textquotedblright{}.) Say it once more.
\end{tcolorbox}
\section[huán]{huán — to give back}

\label{lesson:ZH-C64-huan}



\begin{quote}
Before the new one: say the Mandarin for to do a favour, then the Mandarin for to borrow.
\end{quote}

\subsection*{You'll want to know: huán}
\emph{huán} — \textquotedblleft{}to give back\textquotedblright{}.

\begin{quote}
\emph{wǒ míngtiān huán nǐ} — I'll give it back tomorrow.
\end{quote}

Its characters come later; for now it is a word for the ear.

The same character is read hái in hái xíng; here, giving something back, it is huán.

\begin{cousinweb}
\emph{huán} is inherited from Middle Chinese \emph{hwaen} (Guangyun, in Baxter's transcription), glossed 'to return' and 'to go back'.
\end{cousinweb}

\begin{grammarlens}[title={What you've built}]
One more word for asking politely.
\end{grammarlens}

\subsection*{Your turn}
\begin{itemize}
\raggedright
  \item \emph{Say it:} \emph{huán}
  \item \emph{Say it:} \emph{huán}, once more
  \item \emph{Say it:} \emph{wǒ míngtiān huán nǐ} — I'll give it back tomorrow.
  \item \emph{Recall:} say the Mandarin for to do a favour, then the Mandarin for to borrow, then say \emph{huán} again
\end{itemize}

\begin{tcolorbox}[breakable,colback=teal!4,colframe=teal!35!black,title={Before you move on}]
What does huán mean? (\textquotedblleft{}To give back\textquotedblright{}.) Say it once more.
\end{tcolorbox}
